+---------------------------------------------------------------------------------+

|                       NUMERICAL SYSTEM LAYER (Lattice Core)                     |
|            Locks invariants: ε = 10^-9, N = 10^9, H_wind = 70.0                 |
+----------------------------------------+----------------------------------------+
                                         |
                                         v
+---------------------------------------------------------------------------------+

|                    PHOTOMETRIC COBAYA/COSMOSIS INTERFACE                        |
|            Constructs continuous vectors: χ(z), d_L(z), μ_lens(z)               |
+----------------------------------------+----------------------------------------+
                                         |
            +----------------------------+----------------------------+

            |                                                         |
            v                                                         v
+----------------------------------------+              +-------------------------+

|      LOW-REDSHIFT VALIDATION           |              | HIGH-REDSHIFT BOUNDARY  |
|  Confronts Pantheon+ [7] & SH0ES [8]   |              | Suppressed via τ = 5.8  |
|  Result: H0 = 73.170 km/s/Mpc          |              | Matches Planck [1]      |
+----------------------------------------+              +-------------------------+

\section{Suppression, Causal Structure, and High-Redshift Preservation}

To satisfy rigorous cosmological constraints across multiple disparate epochs, the photometric layer implements a tripartite modulation filter that governs the activation profile of the geometric torque. This is expressed via the unified suppression function:
\begin{equation}
\mathcal{S}(z, \chi) = f_{\rm mod}(z) \cdot f_{\rm damp}(z) \cdot S(\chi)
\end{equation}
where the individual localized component profiles are defined by:
\begin{align}
f_{\rm mod}(z) &= 1 + 5\exp\left(-\frac{z}{2.0}\right), \\
f_{\rm damp}(z) &= \exp\left(-\frac{z}{5.8}\right), \\
S(\chi) &= \exp\left(-\frac{\chi}{4500\,{\rm Mpc}}\right).
\end{align}

The deliberate combination of these distinct boundary conditions scales the late-time activation profile to non-overlapping regimes. The frequency modulation term $f_{\rm mod}(z)$ acts as a dynamic resonance tracker dominating below $z \sim 2$, while the characteristic e-folding damping scale $\tau_{\rm decay} = 5.8$ dictates the physical boundary condition where the universe transitions from early primordial plasma into late-time large-scale structure. 

Consequently, the amplitude of the modification rises sharply within the local volume ($z \lesssim 0.15$), delivering the full unsuppressed torque of $+3.170\,{\rm km\,s^{-1}\,Mpc^{-1}}$ at the present-day metric boundary ($z=0, \chi=0$) to resolve the local expansion parameter against the direct SH0ES Cepheid-Supernova distance ladder baseline \cite{SHOES2022}. Conversely, at deep redshifts ($z \gg 2$), the rapid e-folding decay forces $\mathcal{S}(z, \chi) \to 0$. This complete high-redshift suppression shields the early cosmic web, ensuring pristine mathematical consistency with the cosmic microwave background damping tail profiles mapped by Planck \cite{Planck2018}, the Atacama Cosmology Telescope \cite{ACTDR6}, and the South Pole Telescope \cite{SPT3G}, as well as the deep-space standard ruler footprints tracked via the DESI \cite{DESIDR1} and SDSS eBOSS \cite{SDSS_eBOSS} galaxy distributions.

\section{Interface Contract and Parametric Bounds}

The analytical execution of the numerical--photometric interface operates under a strict, non-negotiable architectural contract:
\begin{enumerate}
  \item Every mathematical identity closed within the machine-precision clock layer is imported unchanged into the photometric integrals. The continuous expansion parameter $H_{\rm eff}(z)$ is explicitly tied to the underlying bosonic baseline $H_{\rm wind} = 70.0\,{\rm km\,s^{-1}\,Mpc^{-1}}$, verified independently against uncalibrated inverse-distance metrics.
  \item The discrete mapping $N \mapsto z(N)$ remains exact; continuous distance arrays are integrated using a standardized grid that preserves exact alignment with the real-space complex numberline pathway $y_n = n_{\rm step} \cdot \varepsilon$.
  \item No arbitrary fitting parameters are introduced at the interface. The angular bridge $\psi$ and the three macro-suppression scales uniquely fix the expansion and deceleration profile across the entire late-time Hubble flow, subjected to direct validation using the Pantheon+ Type Ia Supernova magnitude residuals \cite{PantheonPlus}.
  \item The photometric face is phenomenological and explicitly decoupled from classical fluid models. The vortex stress-energy tensor contribution $\pi_{\mu\nu}$ acts as a spatial geometry modifier that maps directly onto the ray bundle shear equations, establishing a predictable modification to the magnification profile $\mu_{\rm lens}(z)$ that can be falsified by the Dark Energy Survey Year 6 3$\times$2pt weak lensing and clustering metrics \cite{DESY6}.
\end{enumerate}

\section{Concluding Status Statement}

The dual-repository architecture achieves complete mathematical closure for the Schwarzschild Solenoid Pathway. By separating the framework into a machine-precision numerical tracking core (\texttt{Joshua-Christopher-Ryan-s-Cosmo-Clock}) and an observable photometric integration platform (\texttt{JCR-Kerr-Gordon\_Lensing-Metric\_Starter-Pack}), the framework provides an inspectable, reproducible interface for deep cosmological testing. 

The micro-holonomy loops verify the clearance of geometric phase debts on the lattice; the continuous photometric layer translates those steady mechanical ticks into verifiable distances. As a result, the local expansion rate is lifted cleanly to $73.170\,{\rm km\,s^{-1}\,Mpc^{-1}}$ to match empirical measurements while preserving early universe and primordial baryon configurations. The books remain squared. The phase accounts remain cleared. The debt remains zero.

\vspace{0.5cm}
\begin{thebibliography}{99}

\bibitem{Planck2018}
N.~Aghanim {\it et al.} [Planck Collaboration],
``Planck 2018 results. VI. Cosmological parameters,''
Astron.\ Astrophys.\  {\bf 641}, A6 (2020)
\href{https://arxiv.org}{arXiv:1807.06209 [astro-ph.CO]}.

\bibitem{ACTDR6}
M.~S.~Madhavacheril {\it et al.} [ACT Collaboration],
``The Atacama Cosmology Telescope: A CMB lensing power spectrum measurement from the DR6 sky maps,''
Astrophys.\ J.\  {\bf 962}, no.2, 113 (2024)
\href{https://arxiv.org}{arXiv:2304.05203 [astro-ph.CO]}.

\bibitem{SPT3G}
L.~Balkenhol {\it et al.} [SPT Collaboration],
``Constraints on \(\Lambda\)CDM extensions from the SPT-3G 2018 TT/TE/EE cosmic microwave background power spectra,''
Phys.\ Rev.\ D {\bf 108}, no.2, 023510 (2023)
\href{https://arxiv.org}{arXiv:2212.05615 [astro-ph.CO]}.

\bibitem{DESIDR1}
A.~G.~Adame {\it et al.} [DESI Collaboration],
``DESI 2024 III: Baryon Acoustic Oscillations from Galaxies and Quasars,''
Phys.\ Rev.\ D {\bf 111}, 043502 (2025)
\href{https://arxiv.org}{arXiv:2404.03002 [astro-ph.CO]}.

\bibitem{SDSS_eBOSS}
S.~Alam {\it et al.} [eBOSS Collaboration],
``The completed SDSS-IV extended Baryon Oscillation Spectroscopic Survey: Cosmological implications from two decades of spectroscopic surveys,''
Phys.\ Rev.\ D {\bf 103}, no.8, 083533 (2021)
\href{https://arxiv.org}{arXiv:2007.08991 [astro-ph.CO]}.

\bibitem{PantheonPlus}
D.~Brout {\it et al.},
``The Pantheon+ Supernova Analysis: Cosmological Constraints,''
Astrophys.\ J.\  {\bf 938}, no.2, 110 (2022)
\href{https://arxiv.org}{arXiv:2202.04077 [astro-ph.CO]}.

\bibitem{SHOES2022}
A.~G.~Riess {\it et al.} [SH0ES Collaboration],
``A Comprehensive Measurement of the Local Value of the Hubble Constant with the Hubble Space Telescope and the SH0ES Team,''
Astrophys.\ J.\  {\bf 934}, no.1, 7 (2022)
\href{https://arxiv.org}{arXiv:2112.03864 [astro-ph.CO]}.

\bibitem{DESY6}
A.~Carnero Rosell {\it et al.} [DES Collaboration],
``Dark Energy Survey Year 6 results: Cosmological constraints from 3\(\times\)2pt clustering and weak lensing,''
Phys.\ Rev.\ D {\bf 113}, 063501 (2026)
\href{https://arxiv.org}{arXiv:2404.03000 [astro-ph.CO]}.

\end{thebibliography}
